\begin{figure*}[t]
\centering
\figplaceholder[1.8in]{\textwidth}{Fig.~6 --- real-robot tasks and a paired trial}{%
\textbf{Layout:} full width, two rows of six equal panels with 16:9 crops, so that both rows share one panel size.
\textbf{(a) Top row, six equal panels, one per task (T1--T6):} the real photo at the moment \ours grasps or intercepts the target; overlays: 3--4 semi-transparent past poses of the target from the tracker and an arrow for its motion (\textcolor{obscol}{blue}), drawn as in Fig.~\ref{fig:bench_overview}; task name and motion type under each panel (e.g., ``T5 post-impact interception: deflection after the collision'').
\textbf{(b) Bottom row, a paired trial on one task (\fillin{T5 or T6}; not the linear-motion grasp of Fig.~\ref{fig:teaser}):} left three panels \ours, right three panels the strongest baseline, on the same trial key; in each half, the query frame at $\tau_t$, the frame at the execution time $\tapp$ of the first executed command, and contact; overlays: \textcolor{obscol}{blue} tracked past poses and \textcolor{actcol}{purple} commanded gripper path with a marker where the gripper closes; in the baseline row, a \textcolor{failcol}{red} segment from the closing point to the target with its length in mm. Under each half, a timeline strip marking capture, first execution, and contact, annotated with the method's p95 latency from Table~\ref{tab:real}; outcome and criterion above each half.
\textbf{Rules:} real photos, unretouched; overlays from the logged data of the shown trials; same crop and scale within (b).}
\caption{\textbf{Real-robot experiments.} (a) \ours on the six tasks at contact. (b) \ours
and \fillin{the strongest baseline} on one trial key of \fillin{task}: \fillin{e.g., the
baseline acts XX\,ms later and closes XX\,mm behind the target}.}
\label{fig:real_trials}
\end{figure*}